\section{What one algebraic link gives at a full-group fibre}\label{v7:sec:transfer}

Fix a full-group fibre and use \(X,Z,T_X,T_Z\) for its data. Write \(\theta_X,\theta_Z\) for polarizations. At this fibre \(T_X\) and \(T_Z\) are irreducible Hodge structures, and \(\End_{\mathrm{HS}}(T_X)=\End_{\mathrm{HS}}(T_Z)=F\). In particular they have no rational Tate summand and the relevant Lefschetz pairings are nondegenerate on them.

\begin{proposition}[Transfer of the cubic action]\label{v7:prop:transfer}
If an algebraic correspondence \(X^2\to Z\) induces a nonzero map from \(T_X\) to \(T_Z\) after projection to the indicated cohomological components, then the projector onto \(T_X\) and the action of \(F\) on it are algebraic. This criterion holds for the compositum field used at the CM seed.
\end{proposition}
\begin{proof}
Use the algebraic \(H^1\otimes H^1\) projector on \(X^2\) and the projector of Proposition~\ref{v7:prop:proj} on \(Z\). By multiplicity one, the resulting operator \(f\) vanishes off \(T_X\) in \(H^1\otimes H^1\) and is an isomorphism \(T_X\to T_Z\). An algebraic representative of \(\xi\) would give such an operator by multiplication with a power of the polarization on the source and pushforward.

Use the nondegenerate pairings on degree two given by cup product with the appropriate powers of the polarizations. The adjoint \(f^\dagger\) is algebraic: transpose the cycle and use multiplication by these powers and the inverse Lefschetz isomorphism, which is algebraic for abelian varieties \cite{Kle68}. Put \(C=f^\dagger f\). It is zero off \(T_X\) and invertible on \(T_X\). If \(g(t)\) is the minimal polynomial of its restriction to \(T_X\), then \(g(0)\ne0\). The Chinese remainder theorem supplies a polynomial \(P(t)\) with \(P(0)=0\) and \(tP(t)\equiv1\bmod g(t)\). Hence \(P(C)C\) is the projector onto \(T_X\), and
\[
 P(C)\,f^\dagger m_a f
\]
is multiplication by \(a\) on \(T_X\) and zero off it, where \(m_a\) on \(T_Z\) is algebraic by Proposition~\ref{v7:prop:proj}. The identification of the field actions follows from \(f\in F\iota\). If \(f=m_e\iota\), then \(m_{e^{-1}}f=\iota\) is itself algebraic. Applying the inverse Lefschetz map on the source factor to its correspondence class recovers \(\xi\). This uses no nonsquare condition on the polarization or the chosen extension \(K/F\).
\end{proof}

\subsection{From the cubic action to all powers}

We use the following tensor conventions. The polarization gives an alternating form \(\psi\) on \(V\). It and its inverse tensor are algebraic operations: one contracts two \(H^1\) factors by pulling back along a diagonal, multiplying by a polarization power and integrating, and the inverse tensor is a rational multiple of the corresponding mixed Poincar\'e class. Factor permutations and the \(H^1\) K\"unneth projectors are algebraic. Thus a correspondence on \(V\otimes V\) can be treated as an algebraic tensor on four factors, and permutations of its slots, including the partial symplectic transpose, are again algebraic. These assertions may also be obtained by using the algebraic inverse Lefschetz maps of the abelian variety to convert its dual \(H^1\) factors.

\begin{theorem}\label{v7:thm:powers}
At a full-group fibre, an algebraic projector onto \(T_X\) together with its algebraic \(F\)-action implies the rational Hodge conjecture on every \(X^m\). Thus the hypothesis of Proposition~\ref{v7:prop:transfer} implies the conjecture for all such powers.
\end{theorem}
\begin{proof}
Extend coefficients to a splitting field for the representation and the three embeddings of \(F\). Write \(s_i\) for the operator that swaps the two copies of \(V_i\) in \(V\otimes V\) while leaving the other factors fixed, and set \(p_i=(1-s_i)/2\). The three field projectors on \(T_X\), obtained by interpolation in a cubic generator, act on \(V\otimes V\) as
\[
 P_i=(1-p_i)p_jp_\ell,\qquad \{i,j,\ell\}=\{1,2,3\}.
\]
They are algebraic by hypothesis. Let \(\tau\) denote symplectic transpose in the first \(V\) slot, leaving the second slot fixed. The form \(\psi\) is proportional to the product of the three alternating forms. On a two-dimensional factor, partial symplectic transpose sends the swap \(s\) to \(1-s\). Indeed ordinary partial transpose of the swap is the rank-one operator with tensor \(e_1\otimes e_1+e_2\otimes e_2\) on both sides. Conjugation in the first slot by the alternating matrix changes this to the projector times two onto \(e_1\otimes e_2-e_2\otimes e_1\). Therefore \(\tau(p_i)=s_i/2\). Transpose factors independently for this product, giving
\[
 \tau(P_i)=(1-s_i/2)(s_j/2)(s_\ell/2).
\]
The \(s_i\) commute and have eigenvalues \(+1,-1\). The square of the last expression equals \(1/64\) on the \(+1\) eigenspace of \(s_i\) and \(9/64\) on its \(-1\) eigenspace. Hence
\[
 p_i=\frac{64\,\tau(P_i)^2-1}{8}.
\]
This realizes each individual swap \(s_i\) algebraically on the whole \(V\otimes V\). Here the identity is on \(V\otimes V\), enforced by the degree projectors.

On \(V^{\otimes 2d}\), apply a swap to adjacent \(V\) slots to permute the \(i\)-th standard factors independently for each \(i\). The invariants in \((V_i)^{\otimes 2d}\) are spanned by tensors pairing the slots in twos with the alternating form \cite{Weyl39}. In odd degree there are no invariants. Start with \(d\) copies of the algebraic inverse tensor of \(\psi\), which supplies the same pairing in all three colors. Independent permutations give all triples of pairings. Hence every invariant tensor in \(V^{\otimes 2d}\) is algebraic over the splitting field. Passing to exterior cohomology on \(X^m\) is algebraic: choose the slots for each of the \(m\) factors and use diagonal pullback and cup product. The maps span the invariant exterior algebra since antisymmetrization is equivariant in characteristic zero. Rational algebraicity follows by scalar descent.
\end{proof}

\subsection{The two classes on the square}

\begin{proposition}\label{v7:prop:square}
At a full-group fibre, \(\dim_\Q \Hdg^2(X^2)=8\) and products of divisors span a six-dimensional subspace. Any single algebraic class outside that span implies algebraicity of both exceptional directions and the Hodge conjecture on every power of \(X\).
\end{proposition}
\begin{proof}
The primitive part \(U\) of \(\bigwedge^2 V\) has three irreducible complex summands, each with one alternating and two symmetric factors in \(V\otimes V\). The remaining summand, with three alternating factors, is the polarization line. Thus \(\End_{\mathrm{HS}}(\bigwedge^2 V)=\Q\oplus F\). The invariants in \(H^4(X^2)\) in K\"unneth bidegrees \((4,0),(3,1),(2,2),(1,3),(0,4)\) have dimensions \(1,1,4,1,1\). For the middle bidegree use the endomorphism computation. For the others, the trivial representation in \(\bigwedge^4 V\) and the copy of \(V\) in \(\bigwedge^3 V\) each have multiplicity one. To see this last count, the weights of \(V\) are the eight sign triples. Count subsets of size four with sum equal to a specified tuple having \(j\) coordinates equal to 2 and the others 0. The counts for \(j=0,1,2,3\) are \(8,4,2,1\). For size three, count tuples with \(j\) coordinates equal to 3 and the others 1, giving \(4,1,0,0\). Subtracting raised weights in each \(\SL_2\) coordinate gives \(8-3\cdot4+3\cdot2-1=1\) and \(4-3\cdot1=1\).

There are three divisor generators on \(X^2\): the two polarizations \(\theta_1,\theta_2\) and the mixed Poincar\'e divisor \(\delta\). Their six quadratic products are independent. The four nonmiddle bidegrees each receive a nonzero product: for the odd bidegrees this follows by hard Lefschetz, or directly by multiplying the pairing tensor by a polarization. In bidegree \((2,2)\), \(\theta_1\theta_2\) is rank one and \(\delta^2\) induces, up to a scalar, the perfect pairing from the second exterior power. On the primitive space this pairing and the Lefschetz pairing are proportional, since they are invariant for \(\operatorname{Sp}(V,\psi)\), under which the primitive second exterior power is irreducible. The exceptional quotient is thus \(F/\Q\) in the primitive endomorphism space.

Project an algebraic class outside the divisor span to bidegree \((2,2)\) and then to \(U\) on both sides, using degree projectors and subtraction of the polarization line. Conversion to an endomorphism by the algebraic Lefschetz maps gives multiplication by an element of \(F\setminus\Q\). Since \(F\) is cubic, its powers generate the whole field. This already realizes both exceptional directions.

For the assertion about all powers, lift the three spectral projectors of this field action into \(V\otimes V\) by pullback under the addition map and projection to \(H^1\otimes H^1\). Addition sends \(v\wedge w\) to \(v\otimes w-w\otimes v\) in that component; the exterior projection and its normalization are algebraic. These projectors act as \(p_i(1-p_j)(1-p_\ell)\). Their partial symplectic transposes are
\[
 (s_i/2)(1-s_j/2)(1-s_\ell/2).
\]
The second and third factors are strictly positive on each joint eigenspace. Polynomial interpolation taking the value 1 on negative eigenvalues and 0 on positive eigenvalues therefore recovers \(p_i\). The independent swaps then give all invariants by the argument of Theorem~\ref{v7:thm:powers}.
\end{proof}
